%=======================================================================
%  Calibration: financial intermediation block
%  Bibliography keys assumed: bocola2016, gertlerkaradi2011, gertlerkiyotaki2010,
%  hatchondomartinez2009, zettelmeyer2013, neumeyerperri2005.
%  Markers [\textsc{tbc}] flag numbers to be filled from the source named in
%  the adjacent footnote before circulation.
%=======================================================================

\subsection{Financial intermediation}
\label{sec:bankcal}

The intermediary block contributes five parameters: the exit (payout) share $f$,
the divertability parameter $\lambda_X$, the entrant transfer $\omega_X$, the
bankers' discount factor $\beta^{\rm int}$, and, through the bond block, the
coupon decay $\delta_b$ that fixes the duration of the sovereign claim. Of these,
only $f$ and $\beta^{\rm int}$ are set directly. The agency parameters
$(\lambda_X,\omega_X)$ are \emph{not} assigned: they are solved so that the
steady state reproduces two targets --- a leverage ratio $\bar\theta$ and a
steady-state intermediation wedge $\bar\sigma\equiv\bar r^{k}-\bar r$ --- which
is the mapping we describe below. We first fix the objects that the data speak
to directly, and then state precisely which target is not data-identified.

\paragraph{Targets taken from the data.}
The risk-free rate is the euro-area real short rate. Deflating the three-month
interbank rate by headline HICP inflation over the pre-crisis single-currency
sample $1999\text{Q}1$--$2007\text{Q}4$ gives an average of $1.11\%$ per annum,
and we set $\bar r=0.003$ per quarter ($1.20\%$ p.a.).\footnote{FRED series
\texttt{IR3TIB01DEM156N} (three-month interbank rate) and
\texttt{CP0000EZ19M086NEST} (euro-area HICP). We use the pre-crisis window
because the post-2008 average real rate is negative ($-1.19\%$ p.a. over
$2010\text{Q}1$--$2012\text{Q}2$) and would not represent the stationary rate
around which the crisis experiment is conducted.} The bankers' discount factor is
then set to $\beta^{\rm int}=1/(1+\bar r)=0.9970$, so that $\beta^{\rm int}R=1$
at the steady state. This is a normalisation rather than an estimate: bankers in
this model are risk-neutral within the period and $\beta^{\rm int}$ stands in for
a constant stochastic discount factor. It is \emph{not} the household discount
factor, which in the incomplete-markets block satisfies
$\beta_{hh}(1+\bar r)<1$ and is solved separately in the second stage of the
steady state.

The sovereign claim is a Hatchondo--Martinez perpetuity whose coupon decays at
rate $1-\delta_b$ \citep{hatchondomartinez2009}, so its Macaulay duration is
$(1+\bar r)/(\bar r+\delta_b)$ quarters. Setting $\delta_b=0.056$ delivers a
duration of $16.9$ quarters, or $4.2$ years, matching the duration of Greek
marketable central-government debt on the eve of the crisis
[\textsc{tbc}].\footnote{Public Debt Management Agency, \emph{Public Debt
Bulletin}; residual maturity of marketable debt. Duration, not average residual
maturity, is the correct counterpart: it is the elasticity of the bond price to
the discount rate, and it is that elasticity which governs the mark-to-market
loss banks take when risk is priced. The long duration is load-bearing --- at
$\delta_b=0.25$ the repricing of the stock shrinks by roughly a factor of six
and the pass-through mechanism disappears.}

The size of the sovereign position is disciplined by bank exposure rather than by
the public debt-to-GDP ratio: what matters for the transmission mechanism is the
share of intermediary assets exposed to the sovereign, and only the
bank-intermediated portion of the debt appears in the model. We set
$\bar B$ so that domestic sovereign holdings are $7.6\%$ of the domestic bank's
total assets [\textsc{tbc}], which corresponds to $24.5\%$ of annual output in
the calibrated steady state.\footnote{European Banking Authority, 2011 EU-wide
stress test and 2011 capital exercise, sovereign-exposure disclosures for the
Greek banking group. We note that the value currently used, $7.6\%$, is the
Italian banking-sector ratio reported by \citet{bocola2016}; replacing it with
the Greek disclosure is a first-order change, since the exposure share scales the
entire balance-sheet channel.} The distinction matters quantitatively: a reading
of this ratio against GDP rather than against bank assets overstates the exposure
roughly threefold, which is enough to make the fiscal relief from a default
outweigh the bank losses and reverse the sign of the mechanism.

The recovery rate in the feared default event is taken from the realised
restructuring. The March 2012 PSI exchange imposed a face-value reduction of
$53.5\%$; \citet{zettelmeyer2013} estimate the net-present-value haircut at
$59$--$65\%$ depending on the discount rate. We set the recovery rate to
$\mathcal{R}=0.45$, at the conservative end of that range.

\paragraph{The agency friction.}
At the steady state the incentive constraint holds with equality, and the
recursion for the franchise value $\bar\varphi$ collapses onto the closed-form
system
\begin{align}
\bar\varphi&=\frac{\bar\Omega\,(1+\bar r)}{1-\bar\mu},
&
\bar\Omega&=\beta^{\rm int}\bigl[f+(1-f)\bar\varphi\bigr],
&
\bar\mu&=\frac{\bar\Omega\,\bar\sigma\,\bar\theta}{\bar\varphi}.
\label{eq:ss-system}
\end{align}
The system in \eqref{eq:ss-system} is triangular rather than genuinely
simultaneous: $\bar\varphi$ and $\bar\Omega$ cancel from the third equation, so
the multiplier is pinned by the targets alone and the franchise value is then
linear in it. Imposing $\beta^{\rm int}(1+\bar r)=1$,
\begin{align}
\bar\mu&=\frac{\bar\sigma\bar\theta}{1+\bar r+\bar\sigma\bar\theta},
&
\bar\varphi&=\frac{f}{f-\bar\mu},
&
\lambda_X&=\frac{\bar\varphi}{\bar\theta}=\frac{f}{\bar\theta\,(f-\bar\mu)},
&
\omega_X&=\frac{\mathcal{D}}{\bar\theta}-(1-f)\,\bar\sigma,
\label{eq:ss-closedform}
\end{align}
with $\mathcal{D}\equiv1-(1-f)(1+\bar r)=f-\bar r(1-f)$ the discount in the
net-worth accumulation identity. Three features of \eqref{eq:ss-closedform} are
worth recording. First, the steady-state multiplier $\bar\mu$ depends only on the
product $\bar\sigma\bar\theta$ --- the excess return on bank equity --- and is
independent of both $f$ and $\beta^{\rm int}$; the leverage and spread targets
are therefore not separately identified by $\bar\mu$, and it is their product
that the data must discipline. Second, $\mathcal{D}>0$, i.e.\ $f>\bar
r/(1+\bar r)$, is necessary for a stationary net-worth distribution: the exit
share must exceed the risk-free rate, or retained earnings compound without
bound. Third, the calibration is feasible only on a bounded region of target
space,
\begin{equation}
\bar\sigma\,\bar\theta\;<\;\frac{\mathcal{D}}{1-f}
\;=\;\frac{f-\bar r(1-f)}{1-f},
\label{eq:feasible}
\end{equation}
which is the condition $\omega_X>0$: beyond it the entrant transfer required to
sustain the targeted leverage turns negative. Condition \eqref{eq:feasible} is
strictly tighter than the condition $\bar\mu<f$ under which $\bar\varphi$ in
\eqref{eq:ss-closedform} is finite and positive, the two differing by
$\bar r/(1-f)$. At our parameters the binding bound is $\bar\sigma\bar\theta<
3.87\%$ per quarter, equivalently an excess return on bank equity below
$15.5\%$ per annum; the calibrated value $\bar\sigma\bar\theta=0.10\%$ per
quarter lies well inside it.

\paragraph{Internal consistency.}
Two expressions for steady-state net worth must agree. The binding incentive
constraint implies $\bar n^{IC}=\lambda_X\bar{\mathcal{A}}/\bar\varphi$, where
$\bar{\mathcal{A}}$ is the divertable asset base; the accumulation identity
\eqref{eq:nw-lom} implies
\begin{equation}
\bar n^{ACC}
=\frac{1}{\mathcal{D}}
 \sum_{j\in\mathcal{J}}\Bigl[(1-f)\bigl(\bar R_{j}-\bar R\bigr)+\omega_X\Bigr]\bar a_j .
\label{eq:ss-nw}
\end{equation}
Because a single $\lambda_X$ applies to all asset classes
\citep[eq.~3]{bocola2016}, every class earns the same excess return
$\bar R_j-\bar R=\lambda_X\bar\mu/\bar\Omega=\bar\sigma$, so both expressions are
proportional to $\bar{\mathcal{A}}$ and the requirement $\bar n^{IC}=\bar
n^{ACC}$ reduces to a scalar equation in $\bar r^{k}$ that is independent of the
scale \emph{and} the composition of the balance sheet. Solving it in the first
stage of the steady state returns $\bar r^{k}=\bar r+\bar\sigma$ and
$\bar\theta_{ss}=\bar\theta$ exactly. This is the inverse of the calibration map
in \eqref{eq:ss-closedform} rather than an independent restriction, and we use it
as a numerical check: the solved steady state reproduces both targets to machine
precision for any portfolio composition. The single-$\lambda$ assumption is what
makes this work, and it is not innocuous --- allowing divertability to differ
across capital, domestic bonds and foreign bonds re-opens a
portfolio-substitution margin that can render sovereign risk expansionary.

\paragraph{Targets not identified by the data.}
Two targets remain. Leverage is set to $\bar\theta=5$ and the exit share to
$f=0.04$, implying an expected banker horizon of $25$ quarters. Neither has a
clean accounting counterpart: $\bar\theta$ is leverage on the divertable asset
base, not the book assets-to-equity ratio of the banking system, which was
substantially higher over this period; and $f$ governs the payout rate at which
franchise value is dissipated rather than an observable exit frequency. We adopt
the posterior means of \citet{bocola2016}, which are close to the values used by
\citet{gertlerkiyotaki2010}.

The intermediation wedge $\bar\sigma$ is the parameter we can least discipline
directly, and we treat it explicitly as such. In the model's steady state the
default probability is zero, so $\bar\sigma$ is the \emph{entire} steady-state
spread of the sovereign yield over the risk-free rate; any measured spread is
therefore an upper bound on it, since observed spreads also compensate credit
risk. Over $1999\text{Q}1$--$2007\text{Q}4$ the Greek ten-year yield exceeded the
German by $50$ basis points on average.\footnote{FRED series
\texttt{IRLTLT01GRM156N} and \texttt{IRLTLT01DEM156N}. For comparison the average
over $2010\text{Q}1$--$2012\text{Q}2$ is $1{,}248$ basis points, peaking at
$2{,}924$ in February 2012.} We set $\bar\sigma=8$ basis points per annum,
following \citet{bocola2016}, which implies a steady-state excess return on bank
equity of $\bar\sigma\bar\theta=40$ basis points per annum and a barely-binding
constraint, $\bar\mu=0.0010$.

This choice is consequential and we report it plainly. Because
$\bar\sigma\bar\theta$ is the rate at which intermediaries rebuild equity out of
retained earnings, it governs the persistence of any balance-sheet shock. At
$\bar\sigma=200$ basis points --- the value that a naive reading of pre-crisis
sovereign spreads might suggest --- the excess return on equity reaches $10\%$
per annum, an equity loss of the size generated by our risk shock is earned back
within two quarters, and the response of output turns positive from the third
year onward. The small value of $\bar\sigma$ is what makes the transmission
mechanism operate through the level of net worth and the price of the sovereign
claim rather than through a rapidly self-correcting spread. The corollary is that
the steady state sits close to the boundary of the occasionally-binding region:
with $\bar\mu=0.0010$ the constraint binds only marginally on average, and along
the risk experiment it becomes slack after roughly four quarters, so the
sustained contraction is carried by bond prices and net worth rather than by a
persistently elevated credit spread. We report sensitivity to $\bar\sigma$ in
Section~[\textsc{tbc}].

\begin{table}[t]
\centering
\caption{Financial-intermediation parameters}
\label{tab:bankcal}
\begin{tabular}{llll}
\hline\hline
Parameter & Value & Target / moment & Source \\
\hline
\multicolumn{4}{l}{\emph{Set directly}}\\
$\bar r$              & $0.0030$ & real 3M rate, $1.11\%$ p.a., 1999--2007 & FRED, own calculation \\
$\beta^{\rm int}$     & $0.9970$ & $\beta^{\rm int}(1+\bar r)=1$            & normalisation \\
$\delta_b$            & $0.056$  & duration $4.2$ years                     & PDMA [\textsc{tbc}] \\
$\mathcal{R}$         & $0.45$   & PSI haircut $53.5\%$ (face)              & \citet{zettelmeyer2013} \\
$\bar B/\text{assets}$& $0.076$  & bank sovereign exposure                  & EBA 2011 [\textsc{tbc}] \\
$f$                   & $0.04$   & banker horizon, 25 quarters              & \citet{bocola2016} \\
\multicolumn{4}{l}{\emph{Calibration targets}}\\
$\bar\theta$          & $5.00$   & intermediary leverage                    & \citet{bocola2016} \\
$\bar\sigma$          & $8$ bp p.a. & steady-state intermediation wedge     & \citet{bocola2016} \\
\multicolumn{4}{l}{\emph{Solved from \eqref{eq:ss-closedform}}}\\
$\lambda_X$           & $0.2051$ & divertable share                         & implied \\
$\omega_X$            & $0.0072$ & entrant transfer                         & implied \\
$\bar\varphi$         & $1.0257$ & franchise value                          & implied \\
$\bar\mu$             & $0.0010$ & incentive-constraint multiplier          & implied \\
\hline\hline
\end{tabular}
\end{table}

%=======================================================================

\subsection{Interest rates, bond prices, and the return on capital}
\label{sec:ratecal}

Five real rates appear in the model: the deposit rate $\bar r$, the yields on the
two sovereign perpetuities $\bar r^{b}_{D}$ and $\bar r^{b}_{F}$, the return on
bank-financed capital $\bar r^{k}$, and the rate on the working-capital loan
$\bar r^{wc}$ \citep{neumeyerperri2005}. They are not five free parameters. Since
a single divertability parameter $\lambda_X$ applies to every asset class, each
class must deliver the same excess return in order for the intermediary to be
willing to hold it, and the steady-state term structure collapses to
\begin{equation}
\bar r^{k}\;=\;\bar r^{b}_{D}\;=\;\bar r^{b}_{F}\;=\;\bar r^{wc}
\;=\;\bar r+\frac{\lambda_X\bar\mu}{\bar\Omega}\;=\;\bar r+\bar\sigma ,
\qquad
\bar Q_b=\frac{\delta_b}{\bar r+\delta_b+\bar\sigma}.
\label{eq:rate-structure}
\end{equation}
The calibration therefore has \emph{two} degrees of freedom in the rate block ---
the level $\bar r$ and the single wedge $\bar\sigma$ --- and the entire
cross-section of returns is a consequence of them. We record this explicitly
because it bounds what the model can be asked to match: it reproduces one
average spread, not the joint behaviour of sovereign, corporate and deposit
rates.

\paragraph{The deposit rate.}
Deposits are the model's risk-free asset: households hold no other claim, and
the deposit rate paid at $t$ is locked in at $t-1$, so it is predetermined with
respect to the period's shocks. We set $\bar r=0.003$ per quarter, or $1.20\%$
per annum in real terms, against a sample average of $1.11\%$ (Section
\ref{sec:bankcal}). Two features are worth noting. First, the household
discount factors $\beta_D$ and $\beta_F$ are not calibrated to this rate: they
are solved, in the second stage of the steady state, so that household savings
clear the union-wide deposit market at $\bar r$. They are equilibrium objects,
and in the presence of idiosyncratic income risk and a borrowing constraint they
satisfy $\beta(1+\bar r)<1$. Second, deposits are own-good claims remunerated at
national rates, tied together by the real parity condition
$(1+\bar r_{D})=(1+\bar r_{F})\,p'/p$, which is the flexible-price image of a
single union policy rate with national inflation differentials. The cross-border
deposit position that this parity permits is the margin through which national
saving and investment decouple.

\paragraph{The sovereign yield.}
The perpetuity price in \eqref{eq:rate-structure} evaluates to $\bar Q_b=0.9459$
and a yield of $1.28\%$ per annum. Because the priced default probability is zero
at the steady state, this yield contains \emph{no} default premium: the entire
sovereign risk premium studied in the paper is generated by the experiment, not
built into the calibration. The steady-state yield is thus a pure intermediation
object, and its data counterpart is a pre-crisis average taken before sovereign
credit risk was priced in the euro area. Over $1999\text{Q}1$--$2007\text{Q}4$ the
Greek ten-year yield, deflated by Greek HICP inflation, averaged $1.60\%$ per
annum, against $1.28\%$ in the model.\footnote{FRED series
\texttt{IRLTLT01GRM156N} and \texttt{CP0000GRM086NEST}.}

The residual matters for the interpretation of $\bar\sigma$, and we state the
tension rather than conceal it. Two independent readings of the data place the
steady-state sovereign wedge near $50$ basis points per annum: the Greek--German
ten-year differential averaged $50$ basis points over the same window, and the
Greek real ten-year yield exceeded the real three-month rate by $49$ basis
points. Both exceed the calibrated $\bar\sigma=8$ basis points by roughly a
factor of six. Neither measure is a clean counterpart --- the model prices a
$4.2$-year-duration claim with risk-neutral, constant-$\bar\Omega$ intermediaries
and therefore generates no term premium, so the observed differentials conflate
term, credit and intermediation components that \eqref{eq:rate-structure} forces
onto a single number. We nevertheless regard $\bar\sigma$ as the least
well-identified parameter of the block, for the reasons given at the end of
Section~\ref{sec:bankcal}, and report sensitivity to it.

\paragraph{The return on capital.}
Capital is predetermined: the stock producing at $t$ was purchased and priced at
$t-1$ \citep[eq.~6]{bocola2016}, so a shock moves impact output through hours
alone and $\bar r^{k}$ is the return on the vintage the intermediary already
holds. From the firm's first-order condition it is the marginal product net of
depreciation, scaled by the inverse markup,
\begin{equation}
\bar r^{k}=\mathit{mc}\,\alpha\,\frac{\bar Y}{\bar K}-\delta,
\qquad \mathit{mc}=\frac{\varepsilon-1}{\varepsilon},
\label{eq:rk-foc}
\end{equation}
so that $\bar r^{k}$, $\delta$ and the effective capital share jointly pin the
capital--output ratio. This gives a second, independent data check on the level
of the rate block, one that does not run through any spread: at
$\alpha=0.30$, $\varepsilon=6$, $\delta=0.025$ and $\bar r^{k}=0.0032$ the
implied capital--output ratio is $2.22$ in annual terms and the investment share
is $\delta\bar K/\bar Y=22.2\%$, against a Greek average of $24.3\%$ of GDP over
$1999$--$2007$.\footnote{FRED series \texttt{GRCGFCFQDSMEI} and
\texttt{GRCGDPNQDSMEI}, both in current prices.} The gross marginal product,
$11.3\%$ per annum, is the object that must cover depreciation of $10\%$ per
annum; the small \emph{net} return is a consequence of the high quarterly
depreciation rate rather than of the small intermediation wedge.

\paragraph{The return on intermediary equity.}
The rate that governs the dynamics is neither $\bar r$ nor $\bar\sigma$
separately but the return on the intermediary's own funds,
\begin{equation}
\bar r^{n}=\bar r+\bar\theta\,\bar\sigma=1.60\%\ \text{per annum},
\label{eq:roe}
\end{equation}
of which $40$ basis points is the excess over the risk-free rate. This is the
rate at which an intermediary rebuilds equity out of retained earnings, and it,
rather than the spread itself, determines how persistent a balance-sheet shock
is. Its smallness is the reason a mark-to-market loss on the sovereign portfolio
propagates for years in this model instead of being earned back within a few
quarters.

\begin{table}[t]
\centering
\caption{Steady-state real rates: model and data}
\label{tab:rates}
\begin{tabular}{lccl}
\hline\hline
 & Model (\% p.a.) & Data (\% p.a.) & Data counterpart \\
\hline
Deposit / risk-free rate $\bar r$        & $1.20$ & $1.11$ & 3M interbank less HICP, 1999--2007 \\
Sovereign yield $\bar r^{b}_D$           & $1.28$ & $1.60$ & GR 10Y less GR HICP, 1999--2007 \\
Return on capital $\bar r^{k}$           & $1.28$ & ---    & implied $K/Y=2.22$, $I/Y=22.2$ vs $24.3$ \\
Working-capital rate $\bar r^{wc}$       & $1.28$ & [\textsc{tbc}] & ECB MIR, NFC loans less HICP \\
Return on bank equity $\bar r^{n}$       & $1.60$ & [\textsc{tbc}] & bank ROE less risk-free \\
\hline
Intermediation wedge $\bar\sigma$        & $0.08$ & $0.50$ & GR--DE 10Y differential, 1999--2007 \\
Bond price $\bar Q_b$                    & \multicolumn{2}{c}{$0.9459$} & duration $4.2$ years \\
\hline\hline
\end{tabular}

\smallskip
\begin{minipage}{0.92\textwidth}
\footnotesize \emph{Note.} All model rates are real and quarterly, reported at
annual rates. The four asset returns coincide by construction: a single
divertability parameter forces every class to carry the same excess return, so
the model has one spread rather than a term structure. ECB MIR interest-rate
statistics begin in January 2003, which is the binding constraint on the
lending-rate window.
\end{minipage}
\end{table}
